\documentclass[UTF8,a4paper,zihao=-4,fontset=none]{ctexart}

\usepackage{geometry}
\geometry{left=2.8cm,right=2.8cm,top=2.8cm,bottom=2.8cm}

\usepackage{setspace}
\onehalfspacing

\usepackage{booktabs}
\usepackage{longtable}
\usepackage{array}
\usepackage{enumitem}
\usepackage{hyperref}
\usepackage[
  style=gb7714-2015,
  backend=biber,
  gbnamefmt=lowercase,
  autocite=plain
]{biblatex}

\setCJKmainfont[AutoFakeSlant=0.2]{Noto Serif CJK SC}
\setCJKsansfont{Noto Sans CJK SC}
\setCJKmonofont{Noto Sans Mono CJK SC}

% 仅用于论文大标题
\newCJKfontfamily\papertitlefont{Noto Sans CJK SC Medium}

\addbibresource{references.bib}

\hypersetup{
  colorlinks=true,
  linkcolor=black,
  citecolor=black,
  urlcolor=blue,
  pdftitle={从《雅舍小品》到〈還鄉〉：梁实秋台湾时期作品小传},
  pdfauthor={},
  pdfkeywords={梁实秋, 台湾时期, 雅舍小品, 莎士比亚翻译, 作品小传}
}

\ctexset{
  section = {
    format = \large\bfseries,
    name = {,、},
    number = \chinese{section}
  },
  subsection = {
    format = \normalsize\bfseries
  }
}

\title{%
  \papertitlefont\zihao{-2}
  从《雅舍小品》到〈還鄉〉：\\
  梁实秋台湾时期作品小传
}
\author{孙宇辰 125033920373}
\date{2026年7月13日}

\begin{document}

\maketitle

\begin{abstract}
本文以梁实秋 1949 年渡台后至 1987 年的作品序列为对象，以“作品小传”的方式重组生平材料与出版脉络，并结合书目节点、文类转换以及序文、开篇和篇目安排，考察其台湾时期文学身份的形成与重组。《雅舍小品》在台湾结集，使“雅舍”成为此后反复延展的散文标识；莎士比亚译本、辞书和教材的编译则与师范大学任教经历相互交织，显示其文字劳动同时面向文学译介、语言学习和语文教育。莎译工程逐步完成后，旧友记忆、文献编纂、旅行书写与“雅舍”系列续作更为集中；《槐園夢憶》与《梁實秋札記》则呈现悼亡书写与报刊札记在对象和文体功能上的差异。晚年的《雅舍談吃》与《英國文學史》《英國文學選》等著述并行，进一步表明其写作不能被归入“闲适小品”的单一叙述。由此，梁实秋台湾时期的作品关系呈现为散文、翻译、辞书教材、现代文学记忆、悼亡书写与英文学术之间的接续、分岔和回环；遗作〈還鄉〉，则以题名和位置为这一序列留下尾声。
\end{abstract}

\noindent\textbf{关键词：}梁实秋；台湾时期；雅舍小品；莎士比亚翻译；作品小传

\section{引言：以作品为中心的梁实秋台湾时期小传}

若只以“雅舍”散文家的形象概括梁实秋的台湾时期，容易把这一阶段看作《雅舍小品》声名的延续。渡台以后，梁实秋持续出版散文、杂文、回忆文字、辞书教材、翻译作品和英文学术著作，相关书目涵盖小品文、莎士比亚译本、英汉辞书、悼亡散文、饮食散文与文学史编选。\autocite{nmtl_liang_chronology,wenhsun_liang_author} 这些作品并非生平年表上彼此孤立的条目，而是在不同制度环境和生活阶段中承担各自的功能：有的延续“雅舍”的命名，有的面向课堂和语言学习，有的围绕现代文学旧友展开回忆与文献整理，有的进入悼亡书写；到了晚年，日常散文与英文学术著述又在同一书目序列中并行。

本文所考察的作品序列自梁实秋渡台后延续至其逝世前后。1949 年 6 月，梁实秋偕程季淑、梁文蔷自广州赴台湾；同年 11 月，此前见于报刊的三十四篇小品文在臺北结集为《雅舍小品》，由正中書局出版。\autocite{nmtl_liang_chronology,liang_yashe_xiaopin_ncl_scan} 1987 年 12 月，丘彥明主编的多人文集《還鄉: 梁實秋專卷》由聯合文學出版，梁实秋遗作〈還鄉〉收录其中。\cite{1987還鄉} 这两个出版节点为台湾时期小传确立了以作品而非单纯年月组织叙述的首尾：开端不是抵台本身，而是报刊旧篇在新的出版场域中首次结集；末端则不以逝世日期为唯一界限，〈還鄉〉这一遗作题名为晚年叙述留下轻微尾声。

所谓“作品小传”，并非舍弃生平材料，而是调整生平材料与作品叙述之间的关系。生命事件主要为作品的出版与传播、文体变化和写作空间提供传记背景。例如，师范学院和师范大学的任教经历，为考察莎士比亚翻译、英文教材、辞书编纂和语文教育之间的关系提供制度背景；程季淑去世与梁实秋再婚的先后，则为《槐園夢憶》与“四宜軒雜記”的文体差异提供时间坐标；亲人重逢和晚年获奖仅作为晚年书目与接受语境的背景。\autocite{nmtl_liang_chronology,ntnu_liang_house_chronology} 至于遗作〈還鄉〉，则仅以题名及其在该多人文集中的收录位置进入尾声。\cite{1987還鄉} 本文正文采用简体中文；繁体中文文献的题名、责任者及出版项依原版著录，直接引文保留原文字形，其余转述和评论均使用简体中文。

相关资料与既有研究提供了几条相互关联的论述路径。梁实秋作家资料页并列呈现其在散文、翻译、文学批评和英语教育等方面的工作，提示梁实秋的文学身份难以由单一文类涵盖。\autocite{wenhsun_liang_author} 关于莎士比亚翻译，白立平从“贊助”与翻译的关系考察胡適对梁实秋莎译的影响。\autocite{bai2001_patronage_translation} Joubin 则将梁译十四行诗置于全球莎士比亚研究的脉络中考察。\autocite{joubin2023_liang_sonnets} 关于台湾散文场域，劉佳蓉以《文星雜誌》与梁实秋为例，讨论战后散文生产和典律模塑。\autocite{liujiarong2017_wenxing_liang} 关于《雅舍談吃》，陳紫華以北平图像为中心，讨论饮食书写与北平城市记忆的关系。\autocite{chenzihua2012_yashe_tanchi} 上述资料与研究从不同侧面呈现梁实秋的作品脉络，也为本文重新辨析台湾时期作品之间的关系提供了参照。

全文按时间分为五个阶段，各阶段均以作品群为中心。第一节考察《雅舍小品》在台湾的出版如何使“雅舍”成为渡台后可继续使用的文学标识。第二节转向莎士比亚译本、辞书教材与师范大学任教之间的联系，说明这些文字劳动如何同时面向读者与课堂。第三节以 1967--1968 年的莎译出版节点为界，梳理《徐志摩全集》《秋室雜憶》《西雅圖雜記》和《雅舍小品續集》所呈现的现代文学记忆与“雅舍”系列续写。第四节以《槐園夢憶》和《梁實秋札記》为中心，考察悼亡书写、再婚前后的生活变化与写作秩序之间的关系。第五节围绕《白貓王子及其他》《雅舍小品三集》《雅舍談吃》《英國文學史》《英國文學選》《雅舍小品四集》等晚年作品展开，并以遗作〈還鄉〉的题名作为轻微尾声。

\section{《雅舍小品》在台湾：旧篇结集与新身份开端}

《雅舍小品》适合作为梁实秋台湾时期小传的开端，并不因为它是来台后的全新创作，而因为此前发表于报刊的旧篇在新地点首次结集，重新成为辨认作者身份的标识。1949 年 11 月，臺北正中書局出版《雅舍小品》，选收三十四篇小品文。\autocite{nmtl_liang_chronology,liang_yashe_xiaopin_ncl_scan} 序文先把这些小品追溯到“居住在北碚雅舍的時候”\autocite{liang_yashe_xiaopin_ncl_scan}，首篇〈雅舍〉又以“我現在住的『雅舍』”\autocite{liang_yashe_xiaopin_ncl_scan}开出叙述位置。前一句指向北碚旧居，后一句转入正在栖身的空间；两处相接，使“雅舍”同时是记忆中的居所、眼前的住所和台湾出版场域中的书名。台湾时期的开端因此不只是地理迁徙，也是一种既有作者称谓在新出版场域中的再流通。

“雅舍”在这里既是书名和系列名，也牵连着可以继续追索的居住空间与个人记忆。从作品序列看，《雅舍小品》在台湾的出现，成为梁实秋此后可延展的散文框架。后来《雅舍小品續集》于 1972 年出版。\autocite{liang1972_yashe_xiaopin_xuji,nmtl_liang_chronology} 《雅舍小品三集》于 1982 年出版，《雅舍小品四集》于 1986 年出版。\autocite{liang1982_yashe_xiaopin_sanji,liang1986_yashe_xiaopin_siji,nmtl_liang_chronology} 这一系列把“雅舍”从初到台湾的旧篇结集延伸到晚年散文结集。由此看，《雅舍小品》的关键不只在一册书的出版，而在“雅舍”成为台湾时期后续散文反复使用的命名方式。

梁实秋初到台湾后的文字劳动并不限于小品文。1949 年，他主持世界書局《英漢四用字典》的改编并增添新字汇；同年还曾任国立编译馆人文组主任、代理馆长。\autocite{nmtl_liang_chronology} 这使他在台湾的最初位置呈现出两个方向：一方面以《雅舍小品》延续散文家的公众形象，另一方面以辞书、编译和教育工作进入制度化的语文生产。若只从散文风格来理解这一时期，便会遮蔽他在工具书和教育系统中承担的语言整理工作。

1950 年以后，这种位置更加清楚。梁实秋于 1950 年秋辞去国立编译馆职务，转入台湾省立师范学院英语系任教；1951 年又兼任英语系主任。\autocite{nmtl_liang_chronology} 1951 年，他还在《自由青年》发表“文學講話”系列，并翻译《蘇俄的強迫勞工》；1952 年，他为余光中诗集《舟子的悲歌》写书评，并在师院安排下迁入云和街 11 号寓所。\autocite{nmtl_liang_chronology} 这些活动并非彼此孤立：在《雅舍小品》出版之后，梁实秋的文学实践很快延伸到课堂、刊物、书评和居住空间之中。

“文學講話”在《自由青年》连载，使梁实秋初到台湾后的文学论述继续面向青年刊物读者。\autocite{nmtl_liang_chronology} 《舟子的悲歌》出版后，梁实秋又为余光中的诗集撰写序言和书评。\autocite{lu2022_yu_liang_humor} 这两类文字表明，《雅舍小品》之后的梁实秋仍参与面向青年读者的文学论述和书评实践；“雅舍”因而不是封闭的书斋称谓，而是与报刊、课堂、出版社和读者相连的作者位置。

“文學講話”系列所申说的文学观，与同年以“李啟純”署名翻译的《蘇俄的強迫勞工》并置，使文学价值的阐述与政治性译事之间显出一层不易调和的张力。\autocite{nmtl_liang_chronology,wen2007_liang_fanzhuti} 温儒敏将这组文章视为近似“文学概论”的长文，认为它系统发挥了梁实秋一贯的文学观与美学观；其中对普遍人性、理性节制和文学价值的强调，可视为其既有主张在渡台后的延续。\autocite{wen2007_liang_fanzhuti} 白立平与單德興对其后《百獸圖》译本的研究，揭示了反共立场与拒绝文学沦为政治宣传之间的张力；由此回看1951年的译事，“李啟純”这一署名的选择，或许还有更为复杂的背景。\autocite{bai2007_liqichun_baishoutu,shan2014_orwell_taiwan}

1952 年迁入云和街 11 号寓所，使梁实秋在师院系统内获得了相对稳定的居住和工作空间。\autocite{nmtl_liang_chronology,ntnu_liang_house_chronology} 这处寓所也为其后长时段的翻译、辞书编纂和教学写作提供了具体的空间背景。第一阶段的结尾因而落在写作空间的形成，而不是初到台湾的孤独感想象。

1949--1952 年并非单纯的“立足台湾”阶段。《雅舍小品》以报刊旧篇结集的方式进入台湾读者视野，辞书改编和师院任教则使梁实秋进入台湾的语文教育系统。散文家的“雅舍”身份由此同出版物、课堂和日常寓所建立新的联系；这些早期工作也显出后来莎士比亚翻译、英文教材与辞书编纂并行的若干线索。

\section{莎士比亚、辞书与师大：作为文化基础设施的文字劳动}

1953--1966 年间，梁实秋一面从事辞书编纂、教材写作、教学和机构工作，一面持续推进莎译；单册译本与阶段性成果陆续面世，最终的四十册全集形态则尚未完成。\autocite{nmtl_liang_chronology,liang1974_on_translating_shakespeare,shakespeare1968_liang_quanji} 由此观察这一时期，重点不仅在莎译本身，也在这项工程如何与面向读者、课堂和制度的其他文字劳动并行。1953 年，梁实秋应遠東圖書公司之邀着手编纂《最新英漢辭典》，并出版《中外語文文學概要》等作品；1954 年，他出版《莎士比亞的戲劇故事》。\autocite{nmtl_liang_chronology} 从题名和出版位置看，《莎士比亞的戲劇故事》可视为正式莎剧译本大量出版前的导读性工作。《最新英漢辭典》《中外語文文學概要》与《莎士比亞的戲劇故事》文类不同，从文类功能看，分别对应辞书查询、文学知识概述和莎士比亚戏剧导读。

这项工作在台湾时期继续推进，其前史则存在计划倡议与实际承担两个起点。梁实秋后来在〈關於莎士比亞的翻譯〉英译本中回忆，自己实际承担莎译始于 1931 年。\autocite{liang1974_on_translating_shakespeare} 同文所叙的 1930 年，则指向胡適以中基会翻译委员会身份提出、组织莎译计划，并与若干学人共同拟议分工的过程。\autocite{liang1974_on_translating_shakespeare} 因而，1930 年宜视为计划倡议和制度网络的起点，1931 年才是梁实秋自述中的译事承担起点。梁实秋的回忆与英译本编辑注对大陆时期早期译本的具体年份所记不一；可以确认的是，1930 年代中期已有若干梁译莎剧进入商務印書館系统。\autocite{liang1974_on_translating_shakespeare} 由此可见，台湾时期的莎译工作并非从零开始，可以理解为战前未竟计划的继续推进。

师大制度环境进一步扩大了这种文字劳动的范围。1955 年，台湾省立师范学院改制为台湾省立师范大学，梁实秋任文学院院长，并创办英语教学中心。\autocite{nmtl_liang_chronology,ntnu_liang_house_chronology} 1956 年，他创立英语研究所并兼任所长，同时出版《初中最新英文法》《高中最新英文法》等教材。\autocite{nmtl_liang_chronology} 1957 年，《亨利四世上篇》出版。\autocite{nmtl_liang_chronology,liang1974_on_translating_shakespeare} 同年，他又创设国语教学中心。\autocite{nmtl_liang_chronology} 把这些事实并读，梁实秋并不是一个只在书房里与莎士比亚相对的译者。他的译本、辞书、英文法和教学机构，构成了战后台湾语文教育和文学译介的一组实践。

1958 年以后，梁实秋辞去文学院院长、英语研究所所长等行政职，仅任教职；同年出版《談徐志摩》和莎剧译本《冬天的故事》。\autocite{nmtl_liang_chronology} 退出部分行政事务，并不等于退出公共文化劳动。现代文学记忆、莎剧翻译和教学活动在这一年同时出现，已经显示这些文字工作并行展开。1959 年《沉思錄》、1960 年《最新英漢辭典》、1961 年《梁實秋選集》、1962 年《清華八年》、1963 年《秋室雜文》等条目，也表明这一时期仍有散文、回忆、选集和工具书的连续生产。\autocite{nmtl_liang_chronology} 因此，若只把 1953--1966 年看成等待莎译完成的长廊，就会看不见梁实秋在翻译工程旁边持续维护的其他文字通道。

《莎士比亞的戲劇故事》与《亨利四世上篇》的先后出现，使“全集”之前的译介层次变得可见。前者偏向导读和普及，后者进入具体剧作翻译；1964 年，文星書店分两批出版莎剧译本二十册，译介工程已有阶段性成果。\autocite{nmtl_liang_chronology,liang1974_on_translating_shakespeare} 到 1967 年，遠東圖書公司出版《莎士比亞全集》三十七册剧作；1968 年，诗作三册补齐，四十册全集始告完备。\autocite{nmtl_liang_chronology,shakespeare1968_liang_quanji,liang1974_on_translating_shakespeare} 这些节点不宜合并为一个完成叙事：1964 年是阶段性出版，1967 年是剧作部分的大规模完成，1968 年才是全集结构的补足。若把它们混写成一次性成果，便会看不见这项翻译工作持续推进的层次。

单册译本解决的大多是某一剧作的阅读，全集则以成套书目的形式呈现莎士比亚作品的整体规模。梁实秋在台湾完成的不是若干孤立剧本，而是一套具有卷册规模和长期出版过程的中文莎士比亚工程。\autocite{shakespeare1968_liang_quanji,nmtl_liang_chronology} 从作品小传角度看，全集把个人生命时间转化为书籍时间：读者看到的是四十册书的完备形态，背后则是从战前已进入商务系统的早期译本、台湾出版、文星阶段性推出到远东补齐诗作的长时段积累。\autocite{liang1974_on_translating_shakespeare,nmtl_liang_chronology,shakespeare1968_liang_quanji}

辞书和教材同样不应被视为文学事业的附属物。1950 年代至 1960 年代，梁实秋持续参与《最新英漢辭典》、英文法、英文读本和学生英汉辞典等工具书、教材的编写或出版。\autocite{nmtl_liang_chronology} 这些文本未必像小品文那样进入文学鉴赏，却在读者训练和语言学习中发挥作用。所谓“文化基础设施”，指的正是这些文本在语言学习和文学阅读中的使用功能：辞典提供词目和释义，英文法与读本服务语言教学，译本则为中文读者提供进入外国文学的路径。它们文类不同，却共同指向战后台湾教育和出版环境中的语言学习与文学阅读。梁实秋同时从事这些工作，使其台湾时期不只面向文学读者，也面向语言学习者；其“作品”概念因而不能只收窄为抒情和叙事作品，还应包括服务语言学习与文学阅读的文字成果。

莎译工程也不宜只归结为个人毅力。1930 年前后，胡適以中基会翻译委员会的名义提出并组织莎译计划，并通过书信协调参与者。\autocite{liang1974_on_translating_shakespeare} 梁实秋后来把胡適的热心发起视为自己投入莎译的重要条件；胡適也曾建议译本增加注释。\autocite{liang1974_on_translating_shakespeare} 翻译委员会的设置与书信往来，使梁实秋的莎译起步同当时的翻译制度和学术网络相连，可以视为这项长期译事的前史背景。白立平从“贊助”与翻译的关系切入胡適对梁实秋莎译的影响。\autocite{bai2001_patronage_translation} Joubin 关于梁译白话十四行诗的英文研究，则将梁译诗作置于全球莎士比亚研究的问题域。\autocite{joubin2023_liang_sonnets} 胡適的倡议与协调属于莎译前史，台湾阶段的主线仍在梁实秋如何长期推进并完成这项工作。

1966 年，梁实秋自台湾师范大学退休，前后任职十七年。\autocite{nmtl_liang_chronology} 但从作品序列看，退休并没有形成断裂；它更像是莎译全集完成前的制度转换。前文所列的任教、机构建设和辞书教材出版，共同构成一条较长的工作线。沿着这条线观察，梁实秋台湾时期的文学身份不再只有“雅舍主人”一面，也包括译者、教师、辞书编纂者和英语文学知识整理者等位置。

1953--1966 年的作品结构本已越出散文集范围。这一阶段既不是两部散文集之间的空白，也不是莎译全集的单纯等待期；导读、剧作翻译、辞书、教材和机构建设共同构成其间的文字实践。梁实秋在台湾的生命安顿，既借小品文延续原有的散文身份，也通过这些长期而以翻译、工具和教学功能为主的工作，与台湾的教育和出版系统建立持续联系。

\section{完成莎翁之后：旧友记忆与“雅舍”重开}

1967--1968 年是梁实秋台湾时期作品结构中的书目分界。1967 年 8 月 6 日，遠東圖書公司为梁译莎士比亚全集三十七册剧作出版举行庆祝会，约有三百余位学者好友参加。\autocite{nmtl_liang_chronology} 1968 年，三册诗作译本补齐，四十册《莎士比亞全集》出版完备。\autocite{nmtl_liang_chronology,shakespeare1968_liang_quanji,liang1974_on_translating_shakespeare} 与这两个出版节点相邻的书目中，旧友记忆、现代文学文献整理和散文结集更为集中地出现。

莎译完成前后，现代文学记忆已较为集中地进入梁实秋的书目。1967 年，梁实秋出版《談聞一多》；1969 年，他与蔣復璁共同主编六卷本《徐志摩全集》，并发表〈《徐志摩全集》編輯經過〉；同年又出版《秋室雜憶》。\autocite{nmtl_liang_chronology} 莎译工程面向跨语言经典的中文化，《談聞一多》《徐志摩全集》《秋室雜憶》等则转向现代文学人物的保存与回忆。两类工作并不互相排斥：前者把莎士比亚带给中文读者，后者把新月及其周边人物带入回忆和文献整理，共同构成梁实秋晚年文字劳动的不同面向。

这种回忆和编纂不宜简单等同于怀旧。《徐志摩全集》的共同主编工作，使梁实秋进入现代文学文献的整理环节；〈《徐志摩全集》編輯經過〉这一题名本身，也把全集形成过程纳入编者需要交代的对象。\autocite{nmtl_liang_chronology} 与个人散文相比，共同主编一部六卷全集使他进入成套整理和编纂他人作品的位置；与莎士比亚翻译相比，编纂对象又转向中国新文学人物。梁实秋在这一时期同时拥有译者、编者和回忆者三种位置，这使他的台湾时期小传不再只是个人作品的增长，也包括他如何整理他人作品和一段现代文学记忆。

1970 年以后，旅行和家庭空间又进入散文写作。1970 年 4 月，梁实秋赴西雅图探望女儿，8 月返台，后撰成《西雅圖雜記》；1972 年，《西雅圖雜記》出版，同年《雅舍小品續集》由臺北正中書局出版。\autocite{nmtl_liang_chronology,liang1972_yashe_xiaopin_xuji} 就作品位置而言，《西雅圖雜記》把海外探亲后的西雅图经验纳入梁实秋散文系统。《雅舍小品續集》则使“雅舍”这一命名在 1972 年继续用于散文结集，不再只停留在 1949 年的旧篇结集上。

同一阶段还出现《略談中西文化》《實秋雜文》《關於魯迅》等著作条目。\autocite{nmtl_liang_chronology} 从题名和书目位置看，莎士比亚工程完成以后，梁实秋并未只转向私人生活叙述；文学批评、中西文化和现代作家评价仍占有一席之地。第三阶段也并非单纯的散文回潮。所谓“旧友记忆与‘雅舍’重开”，不是翻译结束后的退隐，而是在散文和批评写作中重新安排中西文学经验。

这一时期作品类型的扩展，不宜被理解为从严肃翻译转向轻松散文的简单转换。梁实秋在 1969 年前后整理徐志摩、聞一多和新月相关记忆，1972 年又出版旅行与小品作品，由此将过往的文学交游、当下的家庭移动与长期使用的散文名号纳入同一阶段。\autocite{nmtl_liang_chronology} 劉佳蓉关于《文星雜誌》与梁实秋的研究，把 1950、1960 年代的文学生产和战后散文典律化联系起来考察。\autocite{liujiarong2017_wenxing_liang} 这一研究路径提示，梁实秋在台湾散文场域中的位置不只是个体风格的自然延续，也与刊物、出版及战后散文典律的形成有关。

莎士比亚全集逐步完备之际及其后，现代文学记忆、旅行杂记和“雅舍”系列续写在梁实秋的书目中更为集中。这并不是从“文化工程”退回私人小品，而是文字劳动对象的扩展：前一阶段突出译本、辞书和制度，此后散文、回忆和编辑工作又把已逝友人、海外行踪和“雅舍”旧名号组织进台湾时期的晚年写作。

\section{《槐園夢憶》与“四宜軒”：悼亡、再婚和生活秩序重组}

1974--1976 年间，私人生活的转折与文本形态、写作秩序的变化前后相接，《槐園夢憶》是这一阶段的入口。1974 年 4 月 30 日，程季淑在西雅图因意外受伤不治，5 月 4 日葬于槐园；梁实秋于同年 10 月返台，11 月结识韩菁清，12 月完成并出版《槐園夢憶》。\autocite{nmtl_liang_chronology,liang1974_huaiyuan_mengyi} 这些事实构成作品的传记背景，但《槐園夢憶》并未止于丧偶记录。作品开篇由死亡事实进入悼亡叙述，并为追忆建立文体位置。\autocite{liang_huaiyuan_mengyi_fareast_scan_pages}

《槐園夢憶》的题名不宜只作字面拆解。正文开端题作“悼念故妻程季淑女士”\autocite{liang_huaiyuan_mengyi_fareast_scan_pages}，紧接着交代季淑逝世和葬于西雅图槐园的事实。\autocite{liang_huaiyuan_mengyi_fareast_scan_pages} “槐園”在开篇先是具体的墓园空间；“夢憶”也不是泛泛怀旧，可以理解为在死亡事实之后为追忆取得叙述距离。开篇续段转入悼亡传统和情感语气，出现“聖人忘情”\autocite{liang_huaiyuan_mengyi_fareast_scan_pages}、“喪偶之痛”\autocite{liang_huaiyuan_mengyi_fareast_scan_pages}等语汇。这些语汇把突然丧偶的经验导入悼亡文类，避免文本只停留在事件记录层面；梁实秋台湾时期小传的节奏也因此改变：此前作品多在翻译、课堂、旧友和旅行之间展开，此处则转入与亡妻、墓园和记忆秩序密切相关的悼亡书写。

同年出版的《看雲集》表明，《槐園夢憶》并非孤立地悬置在全部写作之外。\autocite{nmtl_liang_chronology} 不过，《槐園夢憶》使这一阶段的作品序列明显出现悼亡书写。开篇把死亡日期、墓园空间和悼亡语汇相继组织起来，使生活事件经由开篇措辞转化为可被书写、出版和阅读的记忆形式，而不只是悲痛心境的直接表露。\autocite{liang_huaiyuan_mengyi_fareast_scan_pages}

1975 年以后，作品序列中出现“四宜軒”这一新的空间名称，文体也转入报刊札记。梁实秋于 1975 年 5 月 9 日与韩菁清举行婚宴，6 月起在《中華日報》副刊开辟“四宜軒雜記”专栏；该专栏文章后来结集为 1978 年時報文化公司出版的《梁實秋札記》。\autocite{nmtl_liang_chronology,liang1978_zhaji} 该书自序称这六十篇札记“原名『四宜軒雜記』”\autocite{liang1981_zhaji_4th}，并说明其陆续发表于《中華日報》副刊。\autocite{liang1981_zhaji_4th} “四宜軒”首先是专栏空间的命名；《梁實秋札記》作为结集本，又把这一命名及其报刊写作带入书籍形态。它与“槐園”的并置，因而表现为悼亡作品与报刊札记在时间和文体上的相邻，而不是婚姻轶事对作品意义的直接替代。

1975 年《梁實秋自選集》的出版，也可视为同一时期自我整理的一部分。\autocite{nmtl_liang_chronology} 从结集方式看，“自選”可以理解为对既有文字的一次回看和选择。将《槐園夢憶》、“四宜軒雜記”和《梁實秋自選集》并置，1974--1975 年间便同时呈现悼亡作品、报刊专栏和自选本：悼亡文本处理失去，专栏、札记与自选本则延续日常写作。梁实秋台湾时期文学身份的重组，正是在这种断裂与连续的交错中显出轮廓。

《槐園夢憶》与《梁實秋札記》并置，可见两种不同的散文入口。前者在开篇以墓园、亡妻和悼亡语汇组织叙述；后者以自序所称的六十篇读书札记为文类背景，首篇〈親切的風格〉则起于文章与信笔风格的讨论。\autocite{liang_huaiyuan_mengyi_fareast_scan_pages,liang1981_zhaji_4th} 与《槐園夢憶》的悼亡开篇相比，〈親切的風格〉首先把注意力转向读书、文风和日常判断。就两个开篇的对照而言，二者呈现的不是简单的“悲伤之后恢复生活”，而是悼亡叙述与报刊札记在对象和文体功能上的差异。

这一阶段尤其需要把生命事件与文本判断区分开来。1974 年《槐園夢憶》的出版与 1975 年“四宜軒雜記”的开栏前后相接，开篇与自序则分别显示悼亡语汇、空间命名与札记文体的差异。\autocite{nmtl_liang_chronology,liang_huaiyuan_mengyi_fareast_scan_pages,liang1981_zhaji_4th} 生命事件为观察作品在这一时点改换书写对象和文体提供传记背景，却不直接等同于作品意义。

\section{晚年丰收与〈還鄉〉：饮食、英文学术和遗作尾声}

1977--1987 年间，梁实秋台湾时期的多种作品类型较为集中地出现。1977 年，他出版《關於白璧德大師》，并为《新月》翻印本写长序；1978 年，《梁實秋札記》和《梁實秋論文學》出版。\autocite{nmtl_liang_chronology,liang1978_zhaji} 这组书目表明，晚年写作不能只从生活小品进入：白璧德、新月和文学评论延续着思想与文学史线索，“四宜軒雜記”一类专栏文字则在此时结集为书。“闲适”二字不足以涵盖这两类写作的并行。

白璧德相关著作、文学论集与散文结集在晚年书目中并行，说明批评文字并未因“雅舍”系列的延续而消失。札记、文学论述和生活小品分别占据不同的书目位置；这一阶段也并非只有“生活化”的晚年散文，文学批评和读书议论仍是作品序列的一部分。

1980 年出版的《白貓王子及其他》和 1982 年出版的《雅舍小品三集》，使晚年书目继续保留小品与“雅舍”系列线索。\autocite{liang1980_baimao,liang1982_yashe_xiaopin_sanji,nmtl_liang_chronology} 1980 年 6 月，梁实秋在香港与分别多年的儿子梁文骐重逢；1982 年又在西雅图与长女梁文茜会面。\autocite{nmtl_liang_chronology,ntnu_liang_house_chronology} 这些亲人重逢事实构成梁实秋晚年家庭流动的生平背景，与上述作品的出版时间相邻，但不直接解释其中的具体篇章。

1983 年，《永恆的劇場——莎士比亞》与《雅舍雜文》等作品出版。\autocite{nmtl_liang_chronology} 这一书目节点表明，莎士比亚并未随 1968 年全集完备而完全退出梁实秋的作品小传，而是以专题著作的形式再次出现。晚年书目由此既包含“雅舍”散文的继续出版，也保留了早年翻译工程的另一种延续。

1984 年，梁实秋获第九届国家文艺贡献奖，《文訊》第 11 期刊出相关专访。\autocite{nmtl_liang_chronology} 奖项与访谈构成台湾文化界评价其长期文字工作的接受背景；书籍与文章之间的关系，则显出这一阶段多种作品类型的并行。

1985 年集中呈现了晚年写作的不同线索。1 月，《雅舍談吃》由九歌出版社出版。\autocite{nmtl_liang_chronology,liang1985_yashe_tanchi} 8 月，《英國文學史》和《英國文學選》由協志出版公司出版。\autocite{nmtl_liang_chronology,liang1985_yingguo_wenxueshi,liang1985_yingguo_wenxuan} 如果只读《雅舍談吃》，饮食和记忆似乎构成梁实秋晚年最醒目的面向；同年出版的两部英文学术著作则显示，文学史著述和作品编选仍与散文写作并行。两部英文学术著作已于 1979 年完成。\autocite{ntnu_liang_house_chronology} 它们因而更适合放回长期教学和英文学术积累中理解。

《雅舍談吃》的位置，并不在于把梁实秋晚年收束为口腹之乐，而在于它借饮食篇章安放日常趣味与旧都记忆。陳紫華以《雅舍談吃》中的北平图像为题展开研究，把讨论引向饮食与城市记忆的关系。\autocite{chenzihua2012_yashe_tanchi} 目录开端连续列出以食物命名的篇章\autocite{liang_yashe_tanchi_scan_unknown}，使饮食经验进入可连续阅读的篇目秩序。〈汤包〉称其为“故都的独门绝活”\autocite{liang_yashe_tanchi_scan_unknown}，“故都”把菜点叙述同旧都记忆相接；〈火腿〉开头从南北饮食习惯差异入手\autocite{liang_yashe_tanchi_scan_unknown}，又把口味转为地域辨析。篇名与开篇措辞共同表明，《雅舍談吃》的日常性不止于口味描述。与同年出版的《英國文學史》《英國文學選》并看，饮食散文和英文学术构成晚年书目中并行的两类文字工作。

1985 年的《英國文學史》和《英國文學選》也不应被看作《雅舍談吃》的反面。两部书表明，日常散文与文学史著述、作品编选在梁实秋晚年书目中并行。\autocite{liang1985_yingguo_wenxueshi,liang1985_yingguo_wenxuan} 这种并置可以修正“闲适散文家”的单线理解：晚年散文并未脱离长期的翻译、教学、辞书和文学史编选，而是与这些知识劳动共同构成其文字生活。

1986 年，《雅舍小品四集》出版。\autocite{nmtl_liang_chronology,liang1986_yashe_xiaopin_siji} 同年另有《雅舍譯詩》《雅舍懷舊——憶故知》和多种英文读本出版。\autocite{nmtl_liang_chronology} 从 1949 年《雅舍小品》到 1986 年《雅舍小品四集》，这个系列跨过了渡台、莎译、悼亡、札记和饮食书写等多个阶段。\autocite{liang_yashe_xiaopin_ncl_scan,liang1972_yashe_xiaopin_xuji,liang1982_yashe_xiaopin_sanji,liang1986_yashe_xiaopin_siji} 与其把这条时间线解释为预先完成的系列计划，不如说，“雅舍”作为命名在台湾时期被反复使用，成为多个阶段都可辨认的散文标识。

1987 年 12 月，丘彥明主编的多人文集《還鄉: 梁實秋專卷》由聯合文學出版，卷中收录梁实秋遗作〈還鄉〉。\cite{1987還鄉} 〈還鄉〉的题名及其在专卷中的收录位置，为晚年作品序列留下轻微尾声。1949 年的《雅舍小品》把“雅舍”带入台湾，1987 年的〈還鄉〉则在题名层面留下收束意味；两者之间并非一条直线发展的乡愁叙事，而是散文、翻译、辞书、悼亡、饮食和英文学术长期并行之后的首尾照应。

晚年书目呈现出多种文类的并行：饮食散文与英文学术并置，“雅舍”系列与译诗、怀旧文字相邻，遗作〈還鄉〉的题名又构成轻微尾声。这些关系使梁实秋台湾时期文学身份在晚年的继续重组获得具体的作品依据，而不必将其归结为某一个事件或某一部作品。

\section{结语：台湾时期的复合型晚年写作}

梁实秋台湾时期的作品小传可以概括为一条“以文字劳动安顿生命”的作品线索。《雅舍小品》在 1949 年由臺北正中書局出版，使“雅舍”这一散文标识进入新的出版和阅读场域。\autocite{liang_yashe_xiaopin_ncl_scan,nmtl_liang_chronology} 1967 年《莎士比亞全集》三十七册剧作出版，1968 年诗作三册补齐，莎译工程遂在台湾形成可供流通的全集形态。\autocite{shakespeare1968_liang_quanji,nmtl_liang_chronology,liang1974_on_translating_shakespeare} 到 1970 年代中后期，《槐園夢憶》的悼亡开篇与其后的《梁實秋札記》结集先后出现。\autocite{liang_huaiyuan_mengyi_fareast_scan_pages,liang1978_zhaji} 1985 年的《雅舍談吃》把饮食记忆带入晚年散文。\autocite{liang1985_yashe_tanchi,chenzihua2012_yashe_tanchi,liang_yashe_tanchi_scan_unknown} 同年的《英國文學史》《英國文學選》则表明晚年英文学术整理仍在继续。\autocite{liang1985_yingguo_wenxueshi,liang1985_yingguo_wenxuan} 《雅舍小品四集》和收入专卷的遗作〈還鄉〉，则在晚年末端并置出“雅舍”系列与题名尾声。\cite{liang1986_yashe_xiaopin_siji,1987還鄉}

梁实秋台湾时期因而不能只归入“闲适小品”的单一叙述。《雅舍小品》和收入专卷的遗作〈還鄉〉，分别居于这一作品序列的开端与末端。\cite{liang_yashe_xiaopin_ncl_scan,1987還鄉} 其间，莎士比亚全集呈现出翻译工程的长时段结构。\autocite{shakespeare1968_liang_quanji,nmtl_liang_chronology} 《槐園夢憶》《梁實秋札記》把悼亡和札记写作纳入同一晚年序列。\autocite{liang_huaiyuan_mengyi_fareast_scan_pages,liang1981_zhaji_4th} 《英國文學史》《英國文學選》又使这一框架越出散文一线。\autocite{liang1985_yingguo_wenxueshi,liang1985_yingguo_wenxuan} 若从作品而非单个事件出发，梁实秋在台湾的晚年身份呈现为复合结构：他既是“雅舍”散文作者，也是莎士比亚译者、辞书和教材编写者、现代文学记忆的整理者、悼亡书写者、饮食散文作者和英文学术编选者。

从《雅舍小品》的序与首篇、《槐園夢憶》的开篇、《梁實秋札記》的自序和首篇，到《雅舍談吃》的篇目安排与开篇措辞，作品之间的关系由书目次序落实到具体文本。梁实秋台湾时期的小传，由此呈现为作品之间的接续、分岔和回环；晚年文学身份的重新安顿，也在这些作品关系中逐步显形。

\clearpage
\printbibliography[title={参考文献}]

\clearpage
\appendix
\section{梁实秋台湾时期作品年表简表}

{\footnotesize
\begin{longtable}{@{}>{\raggedright\arraybackslash}p{1.15cm}>{\raggedright\arraybackslash}p{5.6cm}>{\raggedright\arraybackslash}p{7.05cm}@{}}
\caption{梁实秋台湾时期作品年表简表\autocite{nmtl_liang_chronology,ntnu_liang_house_chronology,1987還鄉}}\\
\toprule
年份 & 主要作品/出版 & 相关活动或事件 \\
\midrule
\endfirsthead
\toprule
年份 & 主要作品/出版 & 相关活动或事件 \\
\midrule
\endhead
\midrule
\multicolumn{3}{r}{续下页}\\
\endfoot
\bottomrule
\endlastfoot
1949 & 《雅舍小品》；主持《英漢四用字典》改编，增添新字汇 & 6 月偕程季淑、梁文蔷抵台；任国立编译馆人文组主任，后代理馆长 \\
1950 & 继续教学、研究、翻译与创作 & 秋辞国立编译馆职；应聘台湾省立师范学院英语系 \\
1951 & “文學講話”六篇连载；译《蘇俄的強迫勞工》 & 50 岁生日宴友；夏兼任台湾省立师范学院英语系主任 \\
1952 & 〈《舟子的悲歌》〉；〈一個《讀者文摘》的讀者的感想〉 & 台湾省立师范学院拨云和街 11 号寓所 \\
1953 & 译《法國共產黨真相》；合编《中外語文文學概要》 & 应遠東圖書公司邀，着手编纂《最新英漢辭典》 \\
1954 & 《莎士比亞的戲劇故事》；《實秋自選集》；《美國是怎樣的一個國家》；合译《現代戲劇》 & 发表〈獨來獨往——讀蕭繼宗《獨往集》〉 \\
1955 & 英译《錦繡河山》；〈讀《媛珊食譜》〉 & 台湾省立师范学院改制为台湾省立师范大学；任文学院院长，创办英语教学中心 \\
1956 & 译《百獸圖》；主编《初中最新英文法》《高中最新英文法》 & 创立台湾省立师范大学英语研究所并兼任所长 \\
1957 & 《亨利四世上篇》；主编《初中英語複習》《高中英語複習》 & 创设国语教学中心 \\
1958 & 《談徐志摩》；译《冬天的故事》 & 辞文学院院长、英语研究所所长等行政职，仅任教职 \\
1959 & 译《沉思錄》 & — \\
1960 & 英译《寶島臺灣》；主编《最新英漢辭典》 & 赴美国华盛顿大学参加“中美学术合作会议”，并探望女儿、女婿 \\
1961 & 《梁實秋選集》；〈瑪克斯·奧瑞利阿斯——一位羅馬皇帝同時是一位苦修哲學家〉 & 秋，专任台湾省立师范大学英语研究所教授 \\
1962 & 《清華八年》；〈記張自忠將軍〉 & — \\
1963 & 〈憶《新月》〉；“華北視察散記”专辑；《秋室雜文》 & — \\
1964 & 《文學因緣》；文星書店出版莎剧译本二十册 & 发表〈《文學因緣》後記〉 \\
1965 & 主编《最新英語讀本》 & — \\
1966 & 主编《莎士比亞誕辰四百週年紀念集》；〈威爾孫與新劍橋本莎士比亞〉 & 自台湾师范大学退休，服务共十七年 \\
1967 & 《談聞一多》；遠東圖書公司出版三十七册《莎士比亞全集》；主编《學生英漢辭典》 & 8 月 6 日举行梁译莎士比亚全集出版庆祝会 \\
1968 & 补译完成《維諾斯與阿都尼斯》《露克利斯》《十四行詩》；四十册《莎士比亞全集》完备 & 发表〈我在小學〉、〈憶張道藩先生〉、〈憶冰心〉 \\
1969 & 共同主编《徐志摩全集》六卷；〈《徐志摩全集》編輯經過〉；《秋室雜憶》 & — \\
1970 & 《略談中西文化》；《實秋雜文》；《關於魯迅》；主编《袖珍英漢辭典》 & 赴西雅图探望女儿，返台后撰成《西雅圖雜記》 \\
1971 & 《實秋文存》；〈《杜詩研究》序〉 & — \\
1972 & 《西雅圖雜記》；《雅舍小品續集》 & 再赴西雅图，居住女儿梁文蔷家中 \\
1974 & 《看雲集》；《槐園夢憶》 & 程季淑去世并葬于西雅图槐园；梁实秋返台，结识韩菁清 \\
1975 & 《梁實秋自選集》；“四宜軒雜記”专栏 & 往返西雅图处理诉讼；与韩菁清举行婚宴；任大同工学院董事长 \\
1977 & 《關於白璧德大師》；《新月》翻印本长序；〈《新月》前後〉 & — \\
1978 & 《梁實秋札記》；《梁實秋論文學》 & “四宜軒雜記”专栏文章结集为《梁實秋札記》 \\
1979 & 完成《英國文學史》《英國文學選》 & 两书前后历时七年，后于 1985 年 8 月出版 \\
1980 & 《白貓王子及其他》；〈讀《怎樣懂中國戲》〉 & 6 月于香港与分别三十一年的儿子梁文骐重逢 \\
1981 & 〈悼沈宗翰〉；〈清華七十〉 & — \\
1982 & 《雅舍小品三集》；主编《高中英文讀本》 & 师大英语研究所、英语系等举办八十大寿祝寿茶会；获台湾省文艺作家协会资深优良文艺工作者荣誉奖；与长女梁文茜会面 \\
1983 & 《雅舍雜文》；《永恆的劇場——莎士比亞》；〈影響我的幾本書〉 & — \\
1984 & 《看雲集》修订本；《文訊》第 11 期专访〈雅舍主人梁實秋〉 & 获第九届国家文艺贡献奖；预立遗书 \\
1985 & 《雅舍談吃》；《雅舍譯叢》；《雅舍散文》；《英國文學史》；《英國文學選》 & 发表〈漫談《英國文學史》〉、〈教書匠〉 \\
1986 & 《雅舍懷舊——憶故知》；《雅舍小品四集》；《雅舍譯詩》；主编《遠東英文讀本》八册 & 接受季季专访；获《中國時報》文学特别贡献奖 \\
1987 & 《雅舍散文二集》；主编《高農英文讀本》《最新英文讀本》《家事英文讀本》；遗作〈還鄉〉收入丘彥明主编《還鄉: 梁實秋專卷》 & 11 月 3 日逝世；各报刊出追悼专辑；举行文学成就研讨会；设立梁实秋文学奖 \\
\end{longtable}
}

\end{document}
